%%%% THIS Section REMOVED FOR NOW DUE TO LACK OF RELEVANCE --- GOOD INFORMATION IN HERE THOUGH %%%%%%%%

\section{\gai Acceleration}
\label{section:34-literature-gai-acceleration}

This Section reviews current efforts to accelerate \gai to identify what the current areas of focus are and to inform which metrics are likely to make the benchmark a useful discriminative tool.
After identifying a bottleneck, the three options to proceed are \textit{exploit the bottleneck}, \textit{elevate the bottleneck} or \textit{accept the bottleneck}~\cite{Goldratt1990,Goldratt2004}.
Accepting the bottleneck is a business decision that may be made if the costs of exploiting or elevating the bottleneck outweigh the benefits. 

\subsection{Exploitation Techniques}
\textit{Exploiting} a bottleneck involves making changes to the implementation that enable it to better utilize the resources available in the environment. 
Parallelism, Cache blocking, graph-reordering, traversal scheduling, compression, sampling, communication-reduction and new programming paradigms are all \textit{exploitation} techniques.

Parallelizing processing in a single-machine environment can bring more computation power to solve problems. 
Almost all personal computers are multi-core, and most commercial machines are multi-socket, leaving ample latent capacity for parallelism to be exploited. 
Parallel approaches to graph processing are essential to \textit{efficient} graph processing but are hampered by problems with load imbalance, synchronization overheads and non-uniform memory accesses~\cite{Beamer2015}.

Cache blocking aims to divide a graph into a collection of subgraphs that are no larger than the available cache block size. 
The general idea is that by decomposing everything to be small enough to fit in the cache, there should be a significant reduction in irregular memory accesses. 
The efficient and effective decomposition remains a challenge~\cite{Balaji2018, Beamer2017, Ham2016, Zhang2020, Sundaram2015}.

Manipulating how the graph is arranged in memory can yield more efficient execution. 
Graph reordering is a preprocessing technique that aims to optimize the order in which vertices are stored in memory to improve spatial locality\cite{Balaji2018}.
Traversal scheduling is an execution-time activity that aims to pick the optimal order to traverse edges to maximize temporal locality in graph processing~\cite{Mukkara2018}.. 
Both are computationally demanding and can outweigh any benefits gained from a well-ordered graph, particularly in dynamic settings. 

% \nrscomment{cut this Section about compression. you already covered compressed formats in the graph representation Section earlier}
% Compression is useful for balancing the size of a graph in memory with the speed of computation. 
% An adjacency list is more compact than an adjacency matrix is more compact than an incidence matrix~\cite{Kepner2015}, but each level of expansion here has additional information.
% The CSR/CSC format is much more compact than a regular adjacency matrix and is the de facto standard representation for graphs~\cite{Saad2003}.
% Compression beyond the CSR/CSC level has shown to be highly effective in parallel graph processing~\cite{Dhulipala2021}, but the encoding/decoding overheads must be offset by computation gains, and so is generally infeasible for sequential implementations~\cite{Shun2015}.

\nrscomment{this is interesting and does feel like its worth mentioning that GNN stuff uses a very different approach than traditional graph tasks. But i dont think you should have a Section on exploit/elevate techniques at all, so im not sure where this would go. Maybe earlier when discussing the problem categories? it seems like it belongs with explaining the GraphAI problems when you get to GNNs "Oh btw they dont ever operate on the full graph b/c it would be intracTable, they do sampling". it gets lost buried down here, and i want to have known this before i hear about datasets and tasks in benchmarks so i can really understand how to interpret the fact that, for instance, GNN benchmark OGB has graphs with 1.6B edges. Thats way less impressive once you know how GNNs sample. hence i think this bit of info belongs much earlier with explaining GNNs, and maybe giving some similar context to some of the tasks might be worthwhile as well.}
Many machine learning models make aggressive use of sampling to reduce the complexity of computation. 
Embedding methods like deepwalk~\cite{Perozzi2014} and node2vec~\cite{Grover2016} use techniques like random walks to build embeddings reflective of nodes and their neighborhoods. 
Work on GNNs finds that by combining topological and attribute information, a sampling approach can yield results very close to the performance of full-graph analysis~\cite{Hamilton2017}.
Sampling may have some drawbacks for tasks that require more exact information, like GNN-based Structural Analysis tasks based on subgraph isomorphism.

Communication reduction algorithms aim to improve the efficiency of distributed (and shared memory) parallel graph processing by reducing the synchronization overheads between worker nodes using techniques like batching~\cite{Cheng2018}, graph partitioning~\cite{Zhang2018}, representation and algorithmic optimizations~\cite{Azad2021}.

Limited research so far has explored the ability to use approximate or incremental algorithms to reduce processing overheads. 
One exploration of incremental graph partitioning for entity resolution on dynamic graphs illustrates the utility of incremental approaches~\cite{Tauer2019}, and incremental algorithms performed much better on the SAGA-Bench dynamic graph benchmark, based on batch-latency metrics~\cite{Basak2020}.

\nrscomment{having gotten through all that, i think the only bit worth saving is the sampling. the rest of it is general background on graph optimization which is only tangentially relevant, unless it directly pertains to the design of your benchmark (i.e. if youre comparing parallel vs sequential CD, then yes, you need to survey SOTA for each of those) -- but it doesnt feel like thats going to be the case}


\subsection{Elevation Techniques}

\nrscomment{i think you should also cut this Section. after reading it, the only thing i learned that i think might matter to benchmarking is the fact most of the GPU isnt being used most of the time during GNN processing. that makes me want a benchmark that would let that sort of thing be easily discovered, and possibly help find ways to use it more effectively. you can probably kill the rest of this Section and move that one bit up to problem (where is the gap in benchmarking currently, assuming whats out there cant tell you that sort of thing easily) or down to abstract methods motivation (why do we need your empirical approach to designing a benchmark), though i lean towards up rather than down. it could probably support either, but its not RW. your diagram of problem type taxonomy that includes traversal, compute intensive, hybrid, probably goes along with that-- since GNNs are hybrid. that Figure could go in prelim, near Graph Operations, with a cursory overview of graph problem classes, where you sort graphAI tasks by compute intensive, etc.. }

\textit{Elevating} a bottleneck involves upgrading the environment to increase the capacity of the bottleneck.
We discuss elevation as a split of `conventional' techniques like distributing computing, adding memory, upgrading the environment and `hardware acceleration' using GPUs and other accelerators as \textit{elevation} techniques.

\subsubsection{Conventional Elevation Techniques}
Distributed computing was a popular early choice, with the success of map-reduce and other frameworks for distributed computing. 
As graphs became larger and unable to fit in the memory of a single machine, distributed graph processing frameworks like Google's Pregel~\cite{Malewicz2010}, Apache Girafe~\cite{Sakr2016} and Pegasus~\cite{Kang2009} emerged. 
Distributed systems incur a communication overhead in return for accessing potentially greater memory and compute resources~\cite{Shun2013, Beamer2015}.

Increasing the memory available to graph processing systems has shown to be effective, with graphs up to $1$Tb in size being processed in a single shared-memory system~\cite{Dhulipala2021}.
Increasing the available memory requires shared memory parallelism to effectively process the larger graphs that can be held now.
With purely sequential algorithms that scale linearly or worse with the number of edges slowing unacceptably as the size of the graph increases~\cite{Hong2012}.

Beyond adding memory, general upgrades to the environment like implementing higher bandwidth network adapters and adding processor cores can elevate bottlenecks in communication and computation.
In general, these improvements scale with the size of the graph, so only offer a temporary or incremental improvement. 
Adding resources to the environment also increases the costs of acquisition, operation and maintenance, and so the benefits from brute-force elevation need to be factored against these costs. 

\subsubsection{Hardware Acceleration techniques for elevating bottlenecks}
Hardware acceleration involves adding specific hardware to optimize processing. 

\textbf{GPUs.}
GPUs are commonly added to improve dense computing processes like training and inference of neural networks~\cite{Dally2021}.
GPUs use high-speed on-chip memory to accelerate their operations, have significantly more memory bandwidth than a standard CPU architecture and are optimized for dense, regular computations. 
These features allow GPUs to \textit{brute-force} their way through some of the memory-boundedness of CPU-based implementations of algorithms~\cite{Aananthakrishnan2020}.
GPUs are used heavily in the training and acceleration of GNNs~\cite{Wu2021a, Baruah2021, Zhang2023}, and have been used to successfully accelerate traditional graph processing~\cite{Fu2014, Wang2016}.
GPUs are, however, suboptimal for graph processing because of limited on-chip memory, irregular computation patterns, load-balancing issues, unpredicTable execution pathways and data movement overheads.

The data that they are processing needs to fit in the on-chip memory, which is typically much smaller than the main memory of the host system making it infeasible for them to process large graphs on a single node~\cite{Aananthakrishnan2020}, forcing parallel processing and the associated communications overheads. 
The limited memory availability makes it particularly difficult to implement algorithms with large intermediate representations~\cite{Sundaram2015} or that use the GPU densification trick~\cite{Aananthakrishnan2020}.  
The sparsity and resultant irregular memory accesses interfere with the memory coalescing process that is critical to efficient GPU computation~\cite{Aananthakrishnan2020}.
To address memory limitations, GPU graph programming frameworks focus on how to manage the number of `active' vertices during a computation to reduce the demands on the memory~\cite{Wang2016} and leverage sparse representations~\cite{Liu2014,Wu2015}.

Irregular computing patterns degrade the efficiency of GPU processing, with high dependence on the shape of the input graph. 
GPU computing is optimized for dense-regular computations that heavily exploit data parallelism. 
However, graph processing involves sparse-irregular computations that feature amorphous data parallelism and struggle with irregular data accesses/control flow~\cite{Wang2016}.
The impact of input data shape is evident, with more regular degree distributions (like a mesh network) seeing throughput increase linearly as the mean degree increases, but falling sharply as dense, skewed distributions inflict load-balancing overheads~\cite{Wu2015}.

Imbalance in the graph because of a skewed degree distribution exacerbates the unpredictability of load-balancing that comes from not being able to pre-assign the threads and can result in uncoalesced memory patterns, and thread divergence, severely hampering GPU performance relative to dense regular tasks~\cite{Fu2014}. 
Load balancing is the primary challenge when using GPUs to execute graph algorithms~\cite{Wang2016}.
The irregular processing patterns mean that the thread assignments for the GPU can't be known until execution time because of their dependence on the data. 
Thread assignment is a common step in GPU pre-processing and deferring it to execution time because of unpredicTable graph dependencies degrades GPU performance relative to a typically dense, regular workload~\cite{Fu2014}.
In a distributed setting, branches cause thread divergence and synchronization overheads limit thread progress~\cite{Aananthakrishnan2020}.

For the use of an accelerator to be beneficial the efficiency gain in execution time must outweigh the cost of the partitioning and data movement from the main memory to the accelerator and back~\cite{Azad2021}.
Despite all of the research into optimizing GPUs for graph processing, a recent study of GNNs shows that GPUs achieve less than 10\% of peak GPU throughput, less than 50\% of peak memory bandwidth and less than 10\% of the GPU is occupied for 87\% of the time~\cite{Huang2021}.

\textbf{FPGAs and ASICs}
Because GPUs are not optimized for graph processing, many efforts to accelerate \gai focus on the design of \textit{graph hardware accelerators} (GHA). 
Typically these either focus on optimizing only traditional graph processing or for GNNs, need to optimize the hybrid compute patterns of irregular computations for graph traversal and dense computations in neural network training.
The two mechanisms typically employed in GHA design are Field Programmable Gate Arrays (FPGAs) and Application Specific Integrated Circuits (ASICs)~\cite{Zhang2023}.

For traditional graph processors, FPGA and ASIC-based accelerators implement several common optimizations that try to address the identified bottlenecks. 
Popular optimizations include decomposing processing to primitive operations and optimizing those~\cite{Dadu2021, Chen2020}, adopting a sparse matrix representation~\cite{Song2016, Ham2016, Sundaram2015}.
Looking specifically at hardware, distributed processing is improved with optical links to reduce communication overheads~\cite{Song2016, Aananthakrishnan2020} and distributed Global addressing systems~\cite{Dysart2016, Aananthakrishnan2020}.
To reduce the pressure on memory bandwidth approaches look at multithreading~\cite{Aananthakrishnan2020} for data paralleism, using an offload engine~\cite{Aananthakrishnan2020}, optimizing Cache use~\cite{Sundaram2015,Dadu2021} or doing away with it entirely in favor of a Processing-In-Memory or other bespoke architecture~\cite{Song2016, Dadu2021, Dysart2016, Aananthakrishnan2020}.

In general, for GNNs, acceleration methodology uses software-hardware codesign to reduce the processing to primitive steps and then optimize each of those operations in hardware. 
Most GNN accelerators split GNN operations into something that looks like an aggregation phase (graph traversals) and a combination phase (neural network calculations).
FPGAs are used to accelerate the primitive operations and the accelerator uses a hybrid architecture to direct the computations to the appropriate part of the processing environment depending on whether it is the data intensive or the compute-intensive phase~\cite{Kiningham2022, Zhang2023, Yan2020, Lin2023, Zhang2023b, Lin2024}
While each approach varies, the common mechanisms in hardware optimization look like quantization to ease memory bottlenecks~\cite{Zhu2023, Zhong2023}, use of buffers to exploit temporal locality~\cite{Chen2022, Zhang2023a, Chen2023, Zhong2023}, and latency-hiding parallelism using inter-phase pipelining~\cite{Kiningham2022} or techniques like systolic arrays for intra-process parallelism~\cite{Yan2020}.

GNN accelerators are closely tied to the specific implementation that they are accelerating~\cite{Abi2023}. 
As a result, an accelerator optimized for Graph Convolution Networks (GCN) may not be effective on a GraphTransformer model.
The same is true for conventional graph accelerators~\cite{Dadu2021}.
The dependence between algorithm and hardware is core to the ASIC paradigm but is challenging with graph processing because of the impact of input data shape on processing. 
Different GNN models may be better suited depending on the shape of the input graph, and an ASIC approach locks the environment into a particular model. 

\subsection{Summary of \gai Acceleration}

%%%%%%%%%% OLD MATERIAL BELOW HERE %%%%%%%%%%


% \textbf{Graph Pre-Processing}
% Many of the computational problems are not from the graph processing itself, but from generating heuristics and optimizing approximate solutions. 

% Another is to use vertex scheduling to determine the optimal order to visit vertices to maximize spatial locality. 
% vertex scheduling can be offline or online, though online approaches are generally infeasible because of their computational cost~\cite{Mukkara2018}.
% Another optimization that aims to reduce memory bandwidth consumption and latency is `cache-blocking' where the graph is decomposed to subgraphs that can fit into the available cache. 
% Closely related is the idea of \textit{graph-partitioning} which seeks to optimize the division of the graph between worker nodes in a parallel setting so that there are minimal cross-worker dependencies~\cite{Balaji2018}.
% As is discussed in Section~\ref{section:124-general-characteristics}, optimal graph partitioning is $NP$-Complete.
% But sub-optimal partitioning leads to load imbalance and contention between workers for high-degree vertices~\cite{Chen2019}.
% So, heuristic approaches approximate the partitions.
% These approaches can involve computationally expensive eigendecompositions~\cite{Ng2001}, or complex test-and-adjust approaches like the \textit{Kernighan-Lin} algorithm~\cite{Kernighan1970} among others.
% The shape of the graph has a significant impact on the `hardness' of partitioning.
% \textit{Scale-free} graphs with small diameters are very difficult to partition, whereas \textit{mesh} graphs with large diameters are more easily separable because of the reduced interdependence of nodes~\cite{Beamer2015}. 
% Indeed, there is a reasonable assertion that the key to improving the effectiveness of graph processing algorithms may depend on finding algorithms that can efficiently partition graphs and process them as independent subgraphs including streaming and distributed partitioning techniques~\cite{Mccune2015}.





% \textbf{Data and Algorithm Dependence.}
% Parallel algorithms are particularly sensitive to the shape of the input data. 
% Highly skewed degree distributions with small diameters are difficult to partition effectively~\cite{Beamer2015, Wu2015}.
% Algorithms that are \textit{compute intensive} like PageRank which need to process all of the vertices in each iteration suffer from poor locality, and the memory latency cannot be hidden because of the poor communication-computation ratio~\cite{Ahn2015}.  

% \textbf{Parallel abstractions.}
% In general, graph computations are difficult to parallelize using traditional approaches because of their irregular structure and poor communication-computation ratio~\cite{Buluc2011}.
% Factors like skewed degree distributions of graphs, heterogeneity of input graph types (scale-free, random, mesh) and maintaining synchronization make developing correct and efficient algorithms very difficult~\cite{Hong2012, Ediger2013}.
% The proliferation of parallel graph processing frameworks and \textit{Domain Specific Languages} (DSL) that aim to abstract away all the hard-to-implement parallel processing operations and builds in expert optimizations so that the end-user-developer doesn't have to~\cite{Nguyen2013, Dhulipala2017}.
% A common approach is to decompose graph operations into a collection of \textit{primitive} or \textit{atomic} operations, and then combine those to compose the required algorithms for processing. 
% Vertex-centric models~\cite{Mccune2015}, edge-centric models, linear-algebraic models~\cite{Mattson2013, Kepner2016, Buluc2011, Azad2021}, models that apply the \textit{Gather Apply Scatter} (GAS), or the \textit{Frontier} model~\cite{Wang2016} all provide an abstraction to facilitate more efficient parallel graph processing. 
% However, the frameworks are limited, with certain classes of algorithms unable to be represented in their abstractions (e.g. SSSP or MST implementations that use priority queues in early linear-algebraic approaches~\cite{Buluc2011}, bucketing algorithms in \textit{frontier} based models~\cite{Dhulipala2017}), impacts on message passing speed and memory consumption~\cite{Wu2015}, and a general finding that an abstracted library does not perform as \textit{efficently} as a hand-tuned natively implemented algorithm optimized for the system and data it is being run on~\cite{Satish2014}.

% \textbf{Distributed vs. Shared-Memory}
% The tradeoffs between distributed and shared memory parallelism have been discussed extensively above, but the core of the tradeoff is that distributed systems incur a network bottleneck that does not impact shared-memory systems in return for accessing potentially greater memory and compute resources~\cite{Shun2013, Beamer2015}.
% Shared memory systems can exacerbate the already limited memory bandwidth issues as more compute resources compete for the limited memory bandwidth~\cite{Shun2015}, and so does not scale with the number of additional compute nodes~\cite{Ahn2015,Kepner2016}.

% \textbf{Data parallelism vs Amorphous Data Paralellism}.
% Most general-purpose parallel computing frameworks were developed for data parallelism. 
% Data parallelism allows the same operation to be performed on multiple data at the same time to increase efficiency. 
% While some graph processing algorithms have data parallelism, many exhibit a more complex parallelism pattern called “amorphous” data-parallelism~\cite{Lenharth2016}.
% \textit{Amorphous} data parallelism is a more flexible approach that can better handle irregularity in the shape of data, at the cost of being processed less efficiently by general-purpose hardware. 

% \subsubsection{GPU Acceleration}
% The use of GPUs to accelerate computation has seen enormous success in the deep learning community in recent years~\cite{Dally2021}.
% GPUs have significantly more memory bandwidth than a standard CPU architecture and are optimized for dense, regular computations. 
% The increased memory bandwidth allows them to \textit{brute-force} their way through some of the memory-boundedness of CPU-based implementations of algorithms~\cite{Aananthakrishnan2020}, but they have several noTable drawbacks. 

% \textbf{GPU Memory.}
% GPUs use high-speed on-chip memory to accelerate their operations. 
% As a result, the data that they are processing needs to fit in the on-chip memory, which is typically much smaller than the main memory of the host system. 
% The limited memory capacity makes it infeasible for them to process large graphs as a single node~\cite{Aananthakrishnan2020}, meaning that they need to be used as part of a parallel processing framework and manage the associated baggage for parallel processing systems discussed above.
% The limited memory availability makes it particularly difficult to implement algorithms with large intermediate representations, as they consume the available memory and fail~\cite{Sundaram2015}.
% The memory issue is exacerbated by the `GPU Trick' of densifying the connectivity matrix for more efficient computation which can increase memory consumption by orders of magnitude. 
% The sparsity of graphs can also interfere with the memory coalescing process that is critical to efficient GPU computation~\cite{Aananthakrishnan2020}.
% As a result, GPU graph programming frameworks focus on how to manage the number of `active' vertices during a computation to reduce the demands on the memory~\cite{Wang2016}.
% It also highlights the importance of efficient sparse representations, with the CSR representation effectively doubling the size of the graph that can be processed by the GPU compared to a sparse adjacency matrix representation~\cite{Liu2014}, which is why it is the favored data format for all BSP implemenations~\cite{Wu2015}.

% It is also important to note that the host system must copy the data across to the GPU before it can be processed, so for the use of an accelerator to be beneficial the efficiency gain in execution time must outweigh the cost of the partitioning and data movement~\cite{Azad2021}.

% \textbf{GPU Compute.}
% GPU computing is optimized for dense-regular computations that heavily exploit data parallelism. 
% Graph processing involves sparse-irregular computations that feature amorphous data parallelism.
% The impact of input data shape is evident, with more regular degree distributions (like a mesh network) seeing throughput increase linearly as the mean degree increases, but falling sharply as dense, skewed distributions inflict load-balancing overheads~\cite{Wu2015}.
% As a result, GPUs are not optimized for graph processing, struggling with irregular data accesses/control flow~\cite{Wang2016}.
% The Bulk Synchronous Parallel (BSP) paradigm employed by most GPU programming frameworks correlates the number of atomic (primitive) operation invocations with the overall performance of the algorithm--approaches that minimize the work they do tend to perform better overall~\cite{Wu2015}.
% The primary challenge when using GPUs to execute graph algorithms is workload distribution~\cite{Wang2016}.
% The irregular processing patterns mean that the thread assignments for the GPU can't be known until execution time because of their dependence on the data. 
% Thread assignment is a common step in GPU pre-processing and deferring to execution time degrades GPU performance relative to a typical dense, regular workload~\cite{Fu2014}.
% Imbalance in the graph because of a skewed degree distribution exacerbates the unpredictability of load-balancing that comes from not being able to pre-assign the threads and can result in uncoalesced memory patterns, and thread divergence, severely hampering GPU performance relative to dense regular tasks~\cite{Fu2014}. 
% In a distributed setting, branches cause thread divergence and synchronization overheads limit thread progress~\cite{Aananthakrishnan2020}.

